\chapter{测试驱动的重构}

\section{问题入口：代码能跑，为什么还不敢改}

很多项目真正的问题不是没有功能，而是不敢改。一个函数几百行，里面同时读文件、解析、计算、输出。你想修一个小 bug，却不知道会不会破坏其他路径。测试驱动的重构不是先写漂亮架构，而是先把行为固定住。

看一个混在一起的函数：

\begin{lstlisting}[language=C++,caption={难以测试的混合函数}]
void process_file(const std::string& path)
{
    std::ifstream input{path};
    std::unordered_map<std::string, int> counts;
    std::string word;
    while (input >> word) {
        normalize(word);
        ++counts[word];
    }
    for (const auto& [word, count] : counts) {
        std::cout << word << " " << count << '\n';
    }
}
\end{lstlisting}

这段代码难测，因为输入来自文件，输出到终端，核心逻辑藏在中间。重构第一步不是改算法，而是找到可分离的纯逻辑。

\section{先建立外部行为测试}

如果暂时无法拆，可以先用端到端测试固定行为：准备小文件，运行程序，比较输出。这个测试可能慢，但它能保护你第一轮重构。然后逐步抽函数。

第一刀通常是把“清洗单词”抽出来：

\begin{lstlisting}[language=C++,caption={先抽纯函数}]
std::string normalize_word(std::string_view word)
{
    std::string result;
    for (unsigned char ch : word) {
        if (std::isalnum(ch)) {
            result.push_back(static_cast<char>(std::tolower(ch)));
        }
    }
    return result;
}
\end{lstlisting}

这个函数可以直接测试，不需要文件和终端。每抽出一个纯函数，系统就多一个稳定支点。

\section{特征测试和单元测试配合}

重构时常用两类测试：

\begin{itemize}
  \item 特征测试：从外部固定现有行为，哪怕内部很乱。
  \item 单元测试：针对抽出来的小函数验证边界。
\end{itemize}

先用特征测试防止大方向跑偏，再用单元测试支撑细节。不要一开始就追求完美测试结构；面对旧代码，第一目标是建立安全网。

\section{依赖注入让 I/O 可替换}

把文件路径改成输入流，就能测试内存字符串：

\begin{lstlisting}[language=C++,caption={把输入流作为依赖}]
WordCounts count_words_from_stream(std::istream& input)
{
    WordCounts counts;
    std::string word;
    while (input >> word) {
        const std::string normalized = normalize_word(word);
        if (!normalized.empty()) {
            ++counts[normalized];
        }
    }
    return counts;
}
\end{lstlisting}

测试：

\begin{lstlisting}[language=C++,caption={用字符串流测试}]
std::istringstream input{"Cpp cpp, memory"};
const WordCounts counts = count_words_from_stream(input);
assert(counts.at("cpp") == 2);
assert(counts.at("memory") == 1);
\end{lstlisting}

依赖注入不是一定要上框架。把 \code{std::istream&}、\code{std::ostream&}、接口引用或函数对象作为参数，就是最朴素的依赖注入。

\section{重构步骤要小}

一次可靠重构可以按这个节奏：

\begin{enumerate}
  \item 写一个能失败的测试，或固定当前行为。
  \item 抽出一个小函数，不改行为。
  \item 运行测试。
  \item 给小函数补边界测试。
  \item 再抽下一步。
\end{enumerate}

不要一边改结构，一边改需求，一边换算法。失败时你不知道是哪一步引入问题。高质量重构靠小步和反馈，不靠一次性大改。

\section{命名是设计反馈}

如果你抽出的函数很难命名，通常说明边界不清。比如 \code{handle_data}、\code{process_stuff}、\code{do_work} 都是危险信号。好名字应该能说明输入和输出之间的关系：

\begin{itemize}
  \item \code{parse_config_line}：解析一行配置。
  \item \code{count_words}：统计单词。
  \item \code{top_words}：取最高频单词。
  \item \code{validate_registration}：校验注册信息。
\end{itemize}

命名不是表面美化，而是检验函数责任是否单一。

\section{重构后的性能复查}

重构可能改变性能。比如把 \code{std::string_view} 错改成 \code{std::string}，可能增加复制；把一次遍历拆成多次遍历，可能增加成本。每次重构后要看两件事：复杂度有没有变，热路径有没有新增大量分配。

如果重构为了可测试性多了一点常数成本，通常可以接受；如果把线性算法变成平方算法，就不能接受。工程判断要明确写出来。

\begin{exercisebox}
把词频项目的 \code{read_words_from_file} 重构为两层：\code{count_words_from_stream} 和 \code{count_words_in_file}。要求核心统计测试不访问文件系统，文件层只测试打不开文件和正常打开两个路径。
\end{exercisebox}
